\documentclass[11pt,a4paper]{article}

% --- Pacotes Fundamentais ---
\usepackage[utf8]{inputenc}
\usepackage[T1]{fontenc}
\usepackage[brazil]{babel}
\usepackage{lmodern}
\usepackage[top=2.5cm, bottom=2.5cm, left=2.5cm, right=2.5cm]{geometry}
\usepackage{amsmath,amssymb}
\usepackage{graphicx}
\usepackage{booktabs}
\usepackage{listings}
\usepackage{xcolor}
\usepackage{fancyhdr}
\usepackage[most]{tcolorbox}
\usepackage{hyperref}

% --- Configurações de Cores e Links ---
\definecolor{primary}{RGB}{16, 54, 98}      % Azul institucional
\definecolor{secondary}{RGB}{41, 128, 185}  % Azul médio
\definecolor{boxbg}{RGB}{245, 248, 252}     % Fundo para caixas de instrução
\definecolor{boxframe}{RGB}{180, 205, 230}  % Borda das caixas de instrução
\definecolor{codebg}{RGB}{248, 249, 250}

\hypersetup{
    colorlinks=true,
    linkcolor=primary,
    citecolor=secondary,
    urlcolor=secondary
}

% --- Caixa de Instrução para os Alunos ---
\newtcolorbox{instrucao}{
    enhanced,
    colback=boxbg,
    colframe=boxframe,
    arc=3mm,
    boxrule=0.8pt,
    left=4mm, right=4mm, top=3mm, bottom=3mm,
    fonttitle=\bfseries\small\color{primary},
    title={Orientações para esta seção:},
    fontupper=\small\itshape\color{black!80}
}

% --- Configurações para Código-Fonte (Listings) ---
\lstset{
    backgroundcolor=\color{codebg},
    basicstyle=\ttfamily\footnotesize,
    breaklines=true,
    frame=single,
    rulecolor=\color{black!20},
    numbers=left,
    numberstyle=\tiny\color{black!40},
    keywordstyle=\color{primary}\bfseries,
    commentstyle=\color{green!50!black}\itshape,
    stringstyle=\color{orange!80!black}
}
\renewcommand{\lstlistingname}{Listagem}

% --- Cabeçalho e Rodapé ---
\pagestyle{fancy}
\fancyhf{}
\lhead{\footnotesize Projeto em Ciência e Inovação -- 2026/2}
\rhead{\footnotesize Relatório Técnico - Ciclo 3}
\cfoot{\thepage}

% ==============================================================================
% DADOS DO TRABALHO
% ==============================================================================
\title{
    \vspace{-1.0cm}
    \colorbox{yellow}{\textbf{\Large [Título do Projeto / Solução]}}\\[0.3cm] % TODO: discutir o título
    \large Relatório Técnico de Pesquisa e Desenvolvimento (P\&D)
}

\author{
    \textbf{Luis} (Cartão: 000000) -- \texttt{aluno1@inf.ufrgs.br}\\
    \textbf{Manuela} (Cartão: 000000) -- \texttt{aluno2@inf.ufrgs.br}\\
    \textbf{Vítor Feijó} (Cartão: 588403) -- \texttt{vsfeijo@inf.ufrgs.br}\\[0.3cm]
    \normalsize Instituto de Informática -- Universidade Federal do Rio Grande do Sul (UFRGS)\\
    \normalsize Disciplina: Projeto em Ciência e Inovação (2026/2)\\
    \small Repositório do Projeto: \url{https://github.com/vitorsfeijo/PCI-Inovacao-GrupoB}
}

\date{Porto Alegre, \today}

% ==============================================================================
% INÍCIO DO DOCUMENTO
% ==============================================================================
\begin{document}

\maketitle
\thispagestyle{fancy}

\begin{abstract}
\noindent
\textbf{Resumo.} 
\begin{instrucao}
O resumo deve sintetizar o ciclo completo de P\&D (máximo de 250 palavras): (1) O problema real investigado; (2) As hipóteses centrais de solução; (3) A implementação computacional (estruturas de dados e algoritmos); e (4) A metodologia experimental, com os principais resultados e conclusões obtidas.
\end{instrucao}
O planejamento das visitas domiciliares do Agente Comunitário de Saúde (ACS) precisa conciliar a ordem das visitas, as condições da Lei 11.350/2006 e materiais finitos, repostos apenas na Unidade Básica de Saúde. Este relatório apresenta um sistema que modela essas condições em PDDL, de modo que o planejamento automatizado incorpore regras novas sem alterar código. Como escrever cada regra exige dominar PDDL, a hipótese central é que regras descritas em cinco modelos fixos (recurso, preparo, pré-requisito, consumo e exigência por paciente), acrescidos de alternativas (``ou'') e de restrições de horário (turnos da manhã e da noite), podem ser compiladas de forma determinística em PDDL equivalente ao escrito à mão, com menos esforço e sem inviabilizar o planejamento. Implementou-se um compilador que estende um domínio legal travado, instanciação com poda, planejamento por janelas com o Fast Downward e uma interface web para escrever regras e visualizar o roteiro. Na avaliação, com regras e pacientes sintéticos, três regras escritas à mão em PDDL, reescritas como modelos, reproduziram o domínio original em 24 de 24 instâncias; sete regras novas exigiram de 7 a 18 linhas de TOML, contra 13 a 60 linhas de PDDL, e acrescentaram de 1 a 9 minutos a um turno ótimo de 386 minutos. Com as sete regras e 10 pacientes, o roteiro é obtido em cerca de 30 segundos. Também se propõe a opção de um modelo de linguagem, pequeno (SLM) ou grande (LLM), traduzir regras em linguagem natural diretamente para PDDL.
\\[0.3cm]
\textbf{Palavras-chave:} Agentes Comunitários de Saúde; Atenção Primária à Saúde; Planejamento Automatizado; PDDL; Compilação de Regras; Modelos de Linguagem.
\end{abstract}

\vspace{0.3cm}
\hrule
\vspace{0.4cm}

% ------------------------------------------------------------------------------
\section{Contexto e Definição do Problema}

\begin{instrucao}
Apresente o cenário prático real que motivou a investigação e a formalização computacional do problema.
\end{instrucao}

Na Atenção Primária à Saúde do SUS, cada Agente Comunitário de Saúde (ACS) responde por uma microárea e tem na visita domiciliar sua atividade central~\cite{pnab}. Planejar um turno de visitas não se reduz a escolher um percurso: procedimentos como aferir a pressão e medir a glicemia exigem curso técnico, equipamento e assistência de profissional de nível superior (art.~3º, §~4º, da Lei 11.350/2006~\cite{lei11350,lei13595}), e os insumos consumidos só são repostos na Unidade Básica de Saúde (UBS). Cada equipe adota ainda regras próprias, como higienizar as mãos antes dos procedimentos ou usar máscara N95 com pacientes com suspeita de tuberculose. Modeladas como um domínio PDDL~\cite{pddl}, as condições legais são cumpridas pelo planejador sem código específico, mas cada regra nova exige editar o domínio em vários pontos, o que restringe sua manutenção a quem programa.

Como dados clínicos reais são sensíveis, usou-se uma planilha sintética de 10 pacientes fictícios, com marcas como gestante e suspeita de tuberculose; só as coordenadas são reais.

\subsection{Formulação do Problema Computacional}
\begin{itemize}
    \item \textbf{Entradas:} a planilha de pacientes (localização, procedimentos requeridos e marcas), os parâmetros do turno (insumos e janelas de supervisão), a ordem das visitas dada pelo roteirizador e as regras da equipe nos modelos fixos.
    \item \textbf{Restrições:} as condições legais, que as regras não podem afrouxar; a ordem das visitas, que é fixa; e os estoques finitos, repostos só na UBS.
    \item \textbf{Saída:} o plano válido de menor custo em minutos que se consiga obter, apresentado ao ACS como roteiro com o horário e as tarefas de cada visita.
\end{itemize}

% ------------------------------------------------------------------------------
\section{Formulação das Hipóteses de Solução}

\begin{instrucao}
Uma hipótese em computação experimental deve ser mensurável, precisa e passível de validação/falsificação. Defina o que o grupo esperava observar antes de realizar os experimentos.
\end{instrucao}

\subsection{Hipóteses Formuladas}
\begin{itemize}
    \item \textbf{Hipótese Principal ($H_1$, correção):} regras descritas nos cinco modelos fixos e compiladas de forma determinística sobre o domínio legal produzem um domínio equivalente ao escrito à mão em PDDL, em estrutura e em comportamento, em todas as instâncias de teste.
    \item \textbf{Hipótese Secundária ($H_2$, esforço):} descrever uma regra nos modelos exige menos linhas do que as alterações correspondentes no domínio PDDL.
    \item \textbf{Hipótese Secundária ($H_3$, viabilidade):} com as regras ativas, o sistema gera um roteiro válido para um turno de 10 pacientes dentro do limite de 120\,s por busca.
    \item \textbf{Hipótese de Extensão ($H_4$):} alternativas (``ou'') e restrições de horário podem ser incorporadas aos modelos preservando $H_1$ a $H_3$, e um modelo de linguagem consegue traduzir regras descritas em linguagem natural em regras equivalentes às escritas por uma pessoa.
\end{itemize}

\subsection{Métricas de Sucesso}
\begin{itemize}
    \item \textbf{Equivalência ($H_1$):} igualdade de predicados e, em cada ação, de parâmetros, precondições, efeitos e custo; e, em cada instância, mesmo custo ótimo nos dois domínios, com o plano de um válido no outro. Como controle negativo, regras alteradas de propósito devem ser acusadas.
    \item \textbf{Esforço ($H_2$):} linhas escritas no arquivo de regras, sem comentários, contra linhas alteradas no domínio PDDL.
    \item \textbf{Custo e escala ($H_3$):} acréscimo em minutos ao turno ótimo e tempo até o roteiro, variando de 1 a 10 pacientes e de nenhuma a sete regras.
    \item \textbf{Tradução ($H_4$):} proporção de regras geradas pelo modelo de linguagem que compilam e coincidem com um gabarito, desconsiderando os nomes escolhidos.
\end{itemize}

% ------------------------------------------------------------------------------
\section{Desenvolvimento e Implementação da Solução}

\begin{instrucao}
Descreva a realização técnica: como as entidades reais foram mapeadas em estruturas de dados e quais algoritmos foram desenvolvidos.
\end{instrucao}

\subsection{Modelagem Computacional e Arquitetura}
O sistema é um \textit{pipeline}: seleção dos pacientes, roteamento (vizinho mais próximo~\cite{rosenkrantz} com 2-opt~\cite{croes}), compilação das regras, busca do plano, limpeza e geração do roteiro. As regras, em TOML, são aplicadas sobre o domínio legal, que nunca é alterado. Na busca, uma tabela \textit{hash} dá a cada fato um inteiro; um estado é um conjunto imutável desses inteiros, e cada ação instanciada guarda precondições, adições e remoções como conjuntos. Estoques viram níveis discretos encadeados por um predicado de sucessor, mantendo o domínio em STRIPS~\cite{fikes} sem números.

\subsection{Algoritmos e Análise de Custo/Complexidade}
\begin{itemize}
    \item \textbf{Compilação:} cada modelo acrescenta predicados, ações ou precondições, com custo linear no número de regras e de ações, exceto $k$ exigências sobre um mesmo procedimento, que geram $2^k$ versões dele.
    \item \textbf{Instanciação com poda:} em vez de enumerar o produto cartesiano dos parâmetros, liga-os um a um e descarta a combinação assim que uma precondição estática falha. Com sete regras, 4,7 bilhões de combinações viram 9.449 ações.
    \item \textbf{Busca:} A* com $h_{max}$ (admissível)~\cite{bonet} ou Fast Downward~\cite{helmert2006} com LM-cut~\cite{lmcut}, ambos sobre fila de prioridade. Como o planejamento STRIPS é PSPACE-completo~\cite{bylander}, o planejamento por janelas resolve de forma ótima as próximas $k$ visitas e avança uma por vez: custo linear no número de visitas e exponencial só em $k$.
    \item \textbf{Limpeza:} remove ações redundantes~\cite{nakhost}, sem aumentar o custo.
\end{itemize}

\subsection{Tratamento e Pré-processamento dos Dados}
A planilha CSV e os parâmetros JSON são validados na leitura; uma marca sem coluna na planilha é erro. As regras passam por validação em três níveis (campo, regras entre si e conflito com o domínio), com mensagens em português.

% ------------------------------------------------------------------------------
\section{Projeto Experimental e Metodologia de Validação}

\begin{instrucao}
Descreva com rigor o protocolo experimental para assegurar a reprodutibilidade dos resultados.
\end{instrucao}

\subsection{Ambiente Experimental e Variáveis}
\begin{itemize}
    \item \textbf{Ambiente:} Intel Core Ultra 5 245KF (14 núcleos), 16\,GB de RAM, Windows 11, Python 3.14 e Fast Downward pelo pacote \texttt{up-fast-downward}.
    \item \textbf{Variáveis Independentes:} o conjunto de regras ativas (nenhuma, cada uma sozinha ou acumuladas), o número de pacientes (1 a 10), a capacidade dos recursos e a estratégia de busca (A* com $h_{max}$ ou $h_{add}$, Fast Downward com LM-cut ou LAMA~\cite{lama} e janelas com $k$ de 1 a 4).
    \item \textbf{Variáveis Dependentes:} a equivalência estrutural e comportamental, as linhas de TOML e de PDDL alterado, o custo do turno em minutos, o número de ações instanciadas e o tempo de busca.
    \item \textbf{Variáveis de Controle:} a mesma planilha, a mesma rota em todas as execuções, sementes fixas e, nas medições, limite de 120\,s por busca.
\end{itemize}

\subsection{Cenários e Protocolo de Coleta}
Todos os cenários usam dados sintéticos. A equivalência compara, em 24 instâncias geradas com 2 a 4 pacientes e capacidades reduzidas para forçar reposições, o domínio compilado com o escrito à mão; oito alterações propositais nas regras servem de controle negativo. As medições usam sete regras novas (R1 a R7) sobre a planilha de avaliação de 10 pacientes, aplicadas sozinhas, acumuladas uma a uma e com 1 a 10 pacientes. Como planos e custos são determinísticos, cada configuração foi executada uma vez, e os tempos indicam ordem de grandeza. Os planos do Fast Downward são validados no modelo do próprio sistema, e os roteiros das sete regras foram conferidos contra o comportamento esperado de cada uma.

% ------------------------------------------------------------------------------
\section{Resultados Experimentais e Discussão}

\begin{instrucao}
Apresente os dados quantitativos em tabelas e gráficos, confrontando-os diretamente com as hipóteses formuladas na Seção 2.
\end{instrucao}

\subsection{Apresentação dos Dados Coletados}
A Tabela~\ref{tab:regras} mostra o esforço e o impacto de cada regra sozinha, e a Tabela~\ref{tab:acumuladas}, o comportamento do sistema com as regras acumuladas, ambas no turno de 10 pacientes.

\begin{table}[ht]
\centering
\caption{Cada regra sozinha: esforço de escrita e impacto no turno ótimo (sem regras: 386 min)}
\label{tab:regras}
\small
\begin{tabular}{@{}lccccc@{}}
\toprule
\textbf{Regra} & \textbf{TOML} & \textbf{PDDL alterado} & \textbf{Turno ótimo} & \textbf{$\Delta$} & \textbf{Busca ótima} \\
 & (linhas) & (linhas) & (min) & (min) & (s) \\ \midrule
R1 lancetas        & 8  & 27 & 389 & +3 & 0,6 \\
R2 termômetro      & 13 & 36 & 387 & +1 & 1,3 \\
R3 balança         & 7  & 13 & 388 & +2 & 0,3 \\
R4 avental         & 18 & 60 & 395 & +9 & 0,6 \\
R5 N95             & 18 & 60 & 389 & +3 & 1,0 \\
R6 luvas/gestante  & 17 & 52 & 387 & +1 & 1,3 \\
R7 algodão         & 12 & 41 & 389 & +3 & 0,6 \\ \bottomrule
\end{tabular}
\end{table}

\begin{table}[ht]
\centering
\caption{Regras acumuladas: instanciação e busca (custo do turno em minutos e tempo; limite de 120\,s). A última linha acrescenta as três regras usadas no teste de equivalência.}
\label{tab:acumuladas}
\small
\resizebox{\textwidth}{!}{%
\begin{tabular}{@{}lrrccc@{}}
\toprule
\textbf{Regras} & \textbf{Combinações} & \textbf{Ações com poda} & \textbf{A* com $h_{add}$} & \textbf{FD, 1º plano} & \textbf{FD ótimo (LM-cut)} \\ \midrule
nenhuma          & 33.325          & 266    & 427 min, 0,1\,s  & 412 min, 0,1\,s & 386 min, 0,2\,s \\
R1               & 1.777.177       & 1.188  & 444 min, 0,5\,s  & 415 min, 0,1\,s & 389 min, 0,6\,s \\
R1--R4           & 1.780.564       & 1.263  & 537 min, 6,1\,s  & 429 min, 0,1\,s & 401 min, 31\,s \\
R1--R5           & 1.782.377       & 1.305  & 601 min, 34,5\,s & 479 min, 0,1\,s & $>$ 120\,s \\
R1--R7           & 4.676.963.922   & 9.449  & $>$ 120\,s       & 486 min, 0,1\,s & $>$ 120\,s \\
R1--R7 + 3 regras & 565.907.493.560 & 90.545 & $>$ 120\,s      & 550 min, 0,2\,s & $>$ 120\,s \\ \bottomrule
\end{tabular}%
}
\end{table}

\subsection{Validação ou Refutação das Hipóteses}
\begin{itemize}
    \item \textbf{$H_1$ confirmada:} estrutura igual e 24 de 24 instâncias equivalentes; as oito alterações propositais foram acusadas pela comparação estrutural, enquanto a comportamental só as revela quando a instância exercita a regra.
    \item \textbf{$H_2$ confirmada:} 7 a 18 linhas de TOML contra 13 a 60 de PDDL; cada linha escrita evita, em média, 3,1 linhas de PDDL espalhadas por até nove ações.
    \item \textbf{$H_3$ confirmada com o Fast Downward:} o primeiro plano sai em até 0,2\,s em todos os casos, e as janelas ($k=3$) dão 430 min em 16\,s com as sete regras; sem ele, a busca não termina a partir de seis regras.
    \item \textbf{$H_4$:} ainda não avaliada.
    \item \textbf{Complexidade empírica:} as combinações crescem multiplicativamente com os parâmetros, e a poda as torna tratáveis; a busca ótima salta de 31\,s para mais de 120\,s ao passar de quatro para cinco regras, e $k=4$ leva dez vezes mais que $k=3$, como esperado para custo exponencial.
\end{itemize}

\subsection{Discussão Técnica e Limitações Observadas}
O ótimo com muitas regras é desconhecido: os 430 min estão, por estimativa, cerca de 5\% acima dele, e o primeiro plano do Fast Downward fica de 7\% a 45\% acima. A rota é escolhida só pela distância, sem considerar as regras, e voltar da UBS à casa não desfaz preparos como a higienização. Os dados são sintéticos e cada configuração foi executada uma vez.

% ------------------------------------------------------------------------------
\section{Análise de Impacto e Transferência Tecnológica}

\begin{instrucao}
Esta seção avalia o valor gerado além do laboratório. Como a solução computacional impacta a sociedade, a gestão de dados e qual é o caminho para sua aplicação prática?
\end{instrucao}

\subsection{Análise Multidimensional de Impacto}
\begin{itemize}
    \item \textbf{Impacto Social:} beneficia diretamente o ACS e a equipe de Saúde da Família, que mantêm as próprias regras sem depender de quem programa, e indiretamente os pacientes, visitados conforme a lei e os protocolos.
    \item \textbf{Impacto Econômico:} o roteiro substitui o planejamento manual, e o acréscimo de cada regra (1 a 9 minutos) explicita o custo de um protocolo antes de adotá-lo; o ganho frente ao processo manual não foi medido.
    \item \textbf{Ética, Legal e LGPD:} dados de saúde são sensíveis~\cite{lgpd}, e todos os dados usados são fictícios. O sistema roda localmente; a tradução por modelo de linguagem recebe apenas o texto da regra, nunca dados de pacientes, e um SLM local evita enviá-lo a terceiros. Cada ação legal cita seu dispositivo, e o plano é conferível passo a passo. Os critérios de seleção dos pacientes devem ser validados pela equipe, para não excluir ninguém indevidamente.
    \item \textbf{Impacto Ambiental:} indireto, ao evitar idas desnecessárias à UBS, com processamento leve.
\end{itemize}

\subsection{Viabilidade de Implantação e Transferência Tecnológica}
\begin{itemize}
    \item \textbf{Aderência ao Setor:} os parceiros naturais são as secretarias municipais de saúde e o e-SUS Território~\cite{esus}, que registra visitas, mas não planeja rotas nem verifica conformidade normativa.
    \item \textbf{Desafios para Produção:} o protótipo está entre os níveis 3 e 4 de maturidade tecnológica; falta integrar o cadastro do e-SUS, usar distâncias pelas ruas~\cite{osrm} em servidor próprio e operar sem conexão no dispositivo do ACS.
    \item \textbf{Manutenibilidade:} a planilha se atualiza pelo cadastro a cada ciclo de visitas, as regras são revistas quando os protocolos mudam, e o domínio legal, só com a lei.
\end{itemize}

% ------------------------------------------------------------------------------
\section{Conclusões e Lições Aprendidas do Processo de P\&D}

\begin{instrucao}
Sintetize os aprendizados do ciclo de P\&D: o que deu certo, o que não funcionou e quais seriam as próximas investigações.
\end{instrucao}

\subsection{Síntese dos Resultados Técnicos}
Regras descritas em modelos fixos podem ser compiladas em PDDL equivalente ao escrito à mão, com cerca de um terço das linhas e sem afrouxar a lei. O planejamento continua viável com muitas regras, com um planejador de referência e o planejamento por janelas, mas sem garantia de ótimo.

\subsection{O Que Funcionou vs. O Que Não Funcionou}
A tese inicial, de que o planejamento encontraria planos que um programa procedural não encontra, foi refutada: corrigido, o procedural acha os mesmos planos. O que se sustentou foi o custo de manutenção: diante de regras novas, a correção procedural falha, e o domínio PDDL se adapta sem mudar código. Isso motivou o compilador. Também foram reformulados: a instanciação, que não terminava com sete regras; a heurística padrão $h_{add}$, que gerava turnos até 35\% mais longos; e o planejamento só por janelas, que podia esgotar a supervisão antes da última visita.

\subsection{Próximos Passos de Pesquisa e Desenvolvimento}
\begin{itemize}
    \item \textbf{Conflitos entre regras:} detectar, por alcançabilidade relaxada, regras que tornam um paciente impossível de atender ou que não têm efeito.
    \item \textbf{Replanejamento durante o turno:} reagir a paciente ausente, visita urgente ou falta de material a partir do estado atual, medindo tempo de resposta e perda de qualidade.
    \item \textbf{Estudo com dados reais:} com uma UBS parceira, dados anonimizados e aprovação ética, medir o ganho frente ao planejamento manual.
    \item \textbf{Teste com usuários} que não programam, cujo material já está pronto, e a escolha da rota considerando as regras.
\end{itemize}

% ------------------------------------------------------------------------------
\begin{thebibliography}{99}

\bibitem{bonet}
BONET, B.; GEFFNER, H. Planning as heuristic search. \textbf{Artificial Intelligence}, v.~129, n.~1-2, p.~5-33, 2001. DOI: 10.1016/S0004-3702(01)00108-4.

\bibitem{lei11350}
BRASIL. \textbf{Lei nº 11.350, de 5 de outubro de 2006}. Brasília, 2006. Disponível em: \url{https://www.planalto.gov.br/ccivil_03/_ato2004-2006/2006/lei/l11350.htm}. Acesso em: 1 out. 2026.

\bibitem{lei13595}
BRASIL. \textbf{Lei nº 13.595, de 5 de janeiro de 2018}. Altera a Lei nº 11.350, de 5 de outubro de 2006. Brasília, 2018. Disponível em: \url{https://www.planalto.gov.br/ccivil_03/_ato2015-2018/2018/lei/l13595.htm}. Acesso em: 22 set. 2026.

\bibitem{lgpd}
BRASIL. \textbf{Lei nº 13.709, de 14 de agosto de 2018}. Lei Geral de Proteção de Dados Pessoais (LGPD). Brasília, 2018. Disponível em: \url{https://www.planalto.gov.br/ccivil_03/_ato2015-2018/2018/lei/l13709.htm}. Acesso em: 25 nov. 2026.

\bibitem{pnab}
BRASIL. Ministério da Saúde. \textbf{Portaria GM/MS nº 2.436, de 21 de setembro de 2017}. Aprova a Política Nacional de Atenção Básica. Brasília: Diário Oficial da União, 2017. Disponível em: \url{https://www.in.gov.br/materia/-/asset_publisher/Kujrw0TZC2Mb/content/id/19308123}. Acesso em: 22 set. 2026.

\bibitem{esus}
BRASIL. Ministério da Saúde. \textbf{e-SUS Atenção Primária à Saúde: Manual de Uso do Aplicativo e-SUS Território, Versão 3.1}. Brasília: Ministério da Saúde, 2020. Disponível em: \url{https://sisaps.saude.gov.br/esus/upload/docs/manual_utilizacao_aplicativo_esus_aps_territorio.pdf}. Acesso em: 28 set. 2026.

\bibitem{bylander}
BYLANDER, T. The computational complexity of propositional STRIPS planning. \textbf{Artificial Intelligence}, v.~69, n.~1-2, p.~165-204, 1994. DOI: 10.1016/0004-3702(94)90081-7.

\bibitem{croes}
CROES, G. A. A method for solving traveling-salesman problems. \textbf{Operations Research}, v.~6, n.~6, p.~791-812, 1958. DOI: 10.1287/opre.6.6.791.

\bibitem{fikes}
FIKES, R. E.; NILSSON, N. J. STRIPS: a new approach to the application of theorem proving to problem solving. \textbf{Artificial Intelligence}, v.~2, n.~3-4, p.~189-208, 1971. DOI: 10.1016/0004-3702(71)90010-5.

\bibitem{helmert2006}
HELMERT, M. The Fast Downward planning system. \textbf{Journal of Artificial Intelligence Research}, v.~26, p.~191-246, 2006. DOI: 10.1613/jair.1705.

\bibitem{lmcut}
HELMERT, M.; DOMSHLAK, C. Landmarks, critical paths and abstractions: what's the difference anyway? In: \textbf{International Conference on Automated Planning and Scheduling (ICAPS)}, v.~19, p.~162-169, 2009. DOI: 10.1609/icaps.v19i1.13370.

\bibitem{pddl}
McDERMOTT, D. \textit{et al.} \textbf{PDDL: the Planning Domain Definition Language}. Technical Report CVC TR-98-003. New Haven: Yale Center for Computational Vision and Control, 1998.

\bibitem{nakhost}
NAKHOST, H.; MÜLLER, M. Action elimination and plan neighborhood graph search: two algorithms for plan improvement. In: \textbf{International Conference on Automated Planning and Scheduling (ICAPS)}, v.~20, p.~121-128, 2010.

\bibitem{osrm}
PROJECT OSRM. \textbf{Open Source Routing Machine}. Disponível em: \url{https://project-osrm.org}. Acesso em: 28 set. 2026.

\bibitem{lama}
RICHTER, S.; WESTPHAL, M. The LAMA planner: guiding cost-based anytime planning with landmarks. \textbf{Journal of Artificial Intelligence Research}, v.~39, p.~127-177, 2010. DOI: 10.1613/jair.3075.

\bibitem{rosenkrantz}
ROSENKRANTZ, D. J.; STEARNS, R. E.; LEWIS, P. M. An analysis of several heuristics for the traveling salesman problem. \textbf{SIAM Journal on Computing}, v.~6, n.~3, p.~563-581, 1977. DOI: 10.1137/0206041.

\end{thebibliography}

% ------------------------------------------------------------------------------
\section*{Créditos}

Este trabalho foi desenvolvido em dois ciclos da disciplina Projeto em Ciência e Inovação (INF99003) da UFRGS.

\begin{itemize}
    \item \textbf{Ciclo 2:} Gabriel Pieruccini Knopp, Arthur Andrade da Silva e Izadora Candotti de Oliveira desenvolveram o planejador original: a seleção dos pacientes do turno, o roteamento das visitas, os domínios PDDL legal e estendido, o planejador e a planilha de pacientes. Esse código está preservado na pasta \texttt{PCI-Saude-GroupB} do repositório.
    \item \textbf{Ciclo 3:} Luis Henrique Caselani Macedo Junior, Manuela Cristina Zen Dalmonte e Vítor Santana Feijó desenvolveram o compilador de regras, a instanciação com poda, o planejamento por janelas com o Fast Downward, a interface web e a avaliação experimental, além das extensões para alternativas, restrições de horário e tradução de regras por modelos de linguagem.
\end{itemize}

% ------------------------------------------------------------------------------
\appendix
\section{Guia de Reprodutibilidade}

\begin{instrucao}
Forneça os comandos e requisitos necessários para executar o código e reproduzir os experimentos descritos.
\end{instrucao}

\textbf{Requisitos:} Python 3.11 ou superior, sem dependências obrigatórias. O Fast Downward é opcional, mas necessário para reproduzir as colunas ``FD'' da Tabela~\ref{tab:acumuladas} e o experimento de qualidade das janelas; ele é instalado pelo pacote \texttt{up-fast-downward} (versão 1.0.0), que tem versões compiladas para Windows, Linux e macOS. Sem ele, o sistema usa o planejador em Python. No Linux, use um ambiente virtual (\texttt{python3 -m venv .venv}). Os comandos abaixo usam \texttt{python}; no Linux e no macOS, o comando costuma ser \texttt{python3}.

\begin{lstlisting}[language=bash, caption={Instalação e reprodução dos experimentos}]
# 1. Clonar o repositorio
git clone https://github.com/vitorsfeijo/PCI-Inovacao-GrupoB.git
cd PCI-Inovacao-GrupoB/projeto-inovacao-grupoB

# 2. Instalar o Fast Downward (opcional)
python -m pip install -r requirements-opcional.txt

# 3. Testes automaticos (cerca de 1 min)
python acs.py testes

# 4. Experimentos
python acs.py equivalencia   # H1: equivalencia e controle negativo (~3 min)
python acs.py medicoes       # H2 e H3: Tabelas 1 e 2 (~25 min)
python acs.py qualidade      # qualidade das janelas, requer o FD (~40 min)

# 5. Planejar um turno com regras e ver o roteiro
python acs.py planejar exemplos/sinteticas-ciclo2.toml --orcamento 240
python acs.py web            # interface web em http://localhost:8000
\end{lstlisting}

O comando \texttt{medicoes} regrava o arquivo \texttt{exemplos/avaliacao/resultados.md}, de onde vêm as Tabelas~\ref{tab:regras} e~\ref{tab:acumuladas}; a opção \texttt{-{}-rapido} usa limites menores, para conferir o script. As regras R1 a R7 estão em \texttt{exemplos/avaliacao/regras/}, e o \texttt{HOWTORUN.md} detalha todas as opções.

\end{document}
